\chapter{绪论}


\section{研究背景}

% 《普通高中数学课程标准(实验)》指出“高中数学课程应注重提高学生的数学思维能力,这是数学教育的基本目标之一\cite{教育部全日制义务教育数学课程标准}。人们在学习数学和运用数学解决问题时,不断地经历直观感知、观察发现、归纳类比、空间想象抽象概括、符号表示、运算求解、数据处理、演绎证明、反思与建构等思维过程。”数学运算能力是数学基本能力之一，数学运算能力直接决定了学生的数学水平，是学生最基本的数学素养之一, 所以,数学教师必须注重培养学生的数学运算能力。目前我国关于对于西藏高中生这一特殊群体的数学运算能力的研究是寥寥无几的。在经过长期的教学实践探索和改革后,西藏高中生的成绩相对于前几届学生有了明显的提高,大部分学生的数学基础仍然是薄弱的,数学学习成绩是落后的,很多西藏高中生对数学的学习产生了畏难的情绪。为此，笔者尝试剖析西藏高中生的运算能力现状,探索出培养西藏高中生数学运算能力的培养策略。

《普通高中数学课程标准(实验)》提出了“以培养学生的数学思考”为核心的数学教学理念\cite{教育部全日制义务教育数学课程标准}。人类在学习和应用数学的时候，会经历直观感知、观察发现、总结、抽象概括、符号表示、运算求解、数据处理、演绎证明、反思和建构等思考的过程\cite{李明杰2009高中数学课堂学生能力培养的策略研究}。数学计算是数学基础的一种，它是一个人的数学基础，它是一个人的基础。而针对西藏的高中学生，其数学计算能力的问题，在国内尚属空白。西藏高中在多年的教学实践和改革下，与以往相比，取得了较好的成绩，但大多数同学的数学底子还比较差，数学成绩也比较差，许多西藏高中的孩子开始害怕学习。因此，本文试图对西藏高中学生计算技能的状况进行分析，探讨西藏高中学生的数学计算技能的培育对策。



\section{研究意义}

% 教师和学生都知道西藏高中生在数学运算中存在的问题，它是提高数学成绩路上的“拦路虎”，但综合教学进度、学习进度、升学压力等因素，很多教师和学生都忽视了对数学运算能力的培养\cite{董安妮关于西藏高中生数学运算能力现状研究}。数学运算能力贯穿了整个高中数学的学习中，很多学生有清晰的运算思路以及掌握了良好的运算法则，但就是在运算过程上出现问题。当学生出现这种错误时，教师一般会让他们重新计算，可学生往往对自己的错误理解不到位，不知如何修改，几次之后，学生的运算能力就会越来越差，成绩升不上去。可以说想要在数学这一门学科上取得好成绩，就必须要提高数学运算能力。如何采取有效措施培养高中生的数学运算能力是近年来重点研究的问题。本文在以西藏地区西藏高中生、教师、教学环境三个角度分析影响西藏高中生数学运算能力的因素以及培养西藏高中生数学运算能力的方法。

% 师生们都清楚，西藏中学的数学计算是“拦路虎”，但是由于教学进度、学习进度、升学压力等原因，许多老师和同学忽略了对数学运算能力的培养。在中学的教学中，数学的计算是必不可少的，许多同学都有明确的计算方法，也有较强的计算能力。如果是这样的话，老师会要求他们再做一遍，但很多时候，他们对自己的错误并不了解，也不知道该怎么纠正，这样反复多次，结果就是数学能力下降，无法提高分数。可以说，如果你要在这方面有所建树，那么你的数学计算水平就要提升了。目前，针对高中学生的数学计算能力，提出了有针对性的对策。从西藏高中生、教师和教学环境三方面对西藏高中生进行了研究，探讨了对西藏高中生的数学计算能力的影响及其对西藏高中生的数学计算能力的培养。

师生们都清楚，西藏中学的数学计算是“拦路虎”，但是由于教学进度、学习进度、升学压力等原因，许多老师和同学忽略了对数学运算能力的培养\cite{董安妮关于西藏高中生数学运算能力现状研究}。在中学的教学中，数学的计算是必不可少的，许多同学都有明确的计算方法，也有较强的计算能力。如果是这样的话，老师会要求他们再做一遍，但很多时候，他们对自己的错误并不了解，也不知道该怎么纠正，这样反复多次，结果就是数学能力下降，无法提高分数。可以说，如果你要在这方面有所建树，那么你的数学计算水平就要提升了。数学运算能力很大程度上影响学生学习数学的积极性和热情\cite{康恬雨西藏高中生数学运算能力的发展现状与培养策略研究}。如何采取有效措施培养高中生的数学运算能力是近年来重点研究的问题。培养学生的数学计算技能，有利于提高他们的数学思考，使他们的思想素质得到提高，使他们养成一丝不苟、严谨求实的科学态度\cite{龙奕羽素质教育视野中高数教学与课程思政的融合}。目前，针对高中学生的数学运算能力，提出了有针对性的对策。从西藏高中生、教师和教学环境三方面对西藏高中生进行了研究，探讨了对西藏高中生的数学运算能力的影响及其对西藏高中生的数学计算能力的培养。



\section{选题的研究现状}


\subsection{关于数学运算能力的概念内涵}

% 新课标中对运算能力的定义为：“运算能力是指能够根据法则和运算律正确地进行运算的能力。培养运算能力有助于学生理解运算的算理，寻求合理简捷的运算途径解决问题\cite{赵云峰2012培养学生的运算能力初探}。”曹才翰对运算能力的定义为：“是一种综合的能力，是运算技能与逻辑思维能力等的一种独特的结合。

% 陶文中等人对运算能力定义为“运算能力不是简单的加法、减法、乘法、除法的运算过程，而是观察、记忆、推理、表达等功能的从下到上综合发展过程”。

% 李文铭等对运算能力定义：“是计算技能和逻辑思维的综合体现”。计算机技能是指能够记住数学公式和计算规则，并使用概念和特性来计算和推导这些公式和规则，它可以通过合理的过程和严格的推理，正确迅速地得到结果；逻辑思维是指合理运用公式规则进行推理、简化计算、检查和判断计算结果、自我纠正各种计算错误\cite{林瑜提升民族地区西藏高中生数学运算素养的策略研究}。

% 史宁中分别对运算、运算技能、运算能力定义认为“运算是根据已知特定数学概念、规则和原理确定结果的过程；运算技能是通过特定过程和程序的能力；运算能力是一种不仅根据规则、公式等进行计算，而且根据对计算条件的理解来寻找正确的计算方法的能力”。

新课标中对运算能力的定义为：“运算能力是指能够根据法则和运算律正确地进
行运算的能力。”\cite{张劲2021提高运算观察能力} 培养运算能力有助于学生理解运算的算理，寻求合理简捷的运算途径\cite{赵云峰2012培养学生的运算能力初探}。

曹才翰将运算能力界定为：“是一种综合性的能力，它是一种特殊的计算技巧和逻辑思考等。
陶文中等将运算能力界定为：运算能力并非单纯的加法、减法、乘法、除法等操作，而是观察、记忆、推理、表达等功能的全面发展。

李文铭等对运算能力定义：“是计算技能和逻辑思维的综合体现”。计算技能是指能够记住数学公式和计算规则，并使用概念和特性来计算和推导这些公式和规则，它可以通过合理的过程和严格的推理，正确迅速地得到结果；逻辑思维是指合理运用公式规则进行推理、简化计算、检查和判断计算结果、自我纠正各种计算错误\cite{林瑜提升民族地区西藏高中生数学运算素养的策略研究}。

史宁中分别对运算、运算技能、运算能力定义认为“运算是根据已知特定数学概念、规则和原理确定结果的过程；运算技能是通过特定过程和程序的能力；运算能力是一种不仅根据规则、公式等进行计算，而且根据对计算条件的理解来寻找正确的计算方法的能力。”\cite{杨茹冰2021高中生数学运算素养的现状}


\subsection{关于数学运算能力相关的研究}

张莹莹选取杭州和南京各地某学校，运用“数学运算领域”量表，六个维度（加、减、乘、除、填空和大小比较）调查学生数学运算能力。研究发现：男女生数学运算能力无差异；各年级的数学运算能力存在差异，较高年级的运算能力高于较低年级；学生成绩在喜欢数学的程度有明显的差异，喜欢数学的学生比不喜欢数学的学生运算表现好，因此培养学生对数学的兴趣非常重要。

以916位八年级的学生为研究对象，采用问卷调查的方法，从理解操作对象、掌握操作法则、选择操作方法等方面进行了研究。结果表明：八年级的数学计算能力普遍比中等水平强，但男女学生群体中男生的运算基础表现较差；从测试维度分析来看，了解运算对象的表现最好的，而选择运算方法和简便计算上需要更多的关注并努力的得到提高进步。

廉殿勤针对高中第二学段学生的数学运算能力进行调查，以某农村高中为例，运用问卷调查、测试卷、访谈法等方式对高中生七到九年级不同年级学生共400多人为对象，从口算，列竖式计算，简单计算，方程等维度的研究。计算能力的实验结果显示：学生的认知水平和思维水平都随年龄的增加而有所增加，但是操作水平却没有得到改善；在学习过程中，学生在简单的运算法则中，对概念、算理、法则的运用能力较差；男性和女性的计算技能没有差异；不同类型、不同等级的计算技能有明显的正向关系；在错误种类方面，随着年龄的增加，知识和非智力的差错均呈上升趋势。问卷调查显示：尽管随着年龄的增加，大学生的心理发育逐渐趋于完善，对数学计算的重要性有了明确的认知，但是，随着年龄的增加，他们的学习意识和学习意识也会随之增强。